% Pursuing Stacks (cotes 134-2 to 134-8): archive cote 134-8, site page 23 (PDF page 24), typescript page 577.
% typescript. Draft compared: ps7.tex:1068-1126.
% AI-transcribed by Claude Opus 5.5 (claude-opus-5-5, adaptive thinking, effort high), Message Batches API, batch ps4 (msgbatch_012Aoc1NuskATPsyBeKSYK7J), 2026-11-12; not yet checked by eye.
\scanpage{134-8/23}
\mspage{577}

Thus we see that when working with the notion of properness of
maps, even for such simple gadgets as finite ordered sets, which \xout{shoul}
should be viewed as "proper" (or "compact") objects by themselves,
this notion is far from being a wholly trivial one - for instance it
does not hold true that any map between such "proper" objects is again
"proper". This now seems to me just a mathematical "fact of life",
which we may not disregard when working with finite ordered sets say,
in view of expressing in algebraic terms some standard operations \unsure{in|on}
the cohomology theory of sheaves on topological spaces, endowed with
suitable conical stratifications. The fault surely is not with the
notion of conical stratification itself, which may be felt by some
as being ad hoc, awkward and what not. I know the notion is just
right - but at any rate, even when working with strata which are
perfect topological cells (so that nobody could possibly object to
them), exactly the same facts of life are there - \inserted{\deleted{\illegible{1}}}\hand{(}not every map \xout{of}
between such cellular decompositions, mapping cells \uline{onto} cells, will
fit into a"combinatorial"description,when it comes to describing the
standard operations of the cohomology of sheaves, for sheaves compatib-%
le with the stratifications ...

To sum up, it seems to me that definitely, there is a very rich
experimental mat\xout{r}erial available already at present, in order to come
to a feeling of what duality is like in the context of (Cat), and \xout{for}
for developping some of the basic intuitions needed for working out,
hopefully\xout{\unsure{,}}, "the" full-fledged duality theory which should hold in
(Cat). Before leaving this digression, I would like to point out still
one \xout{more} other property connected with the \xout{notion} intuitions around
"properness" - more generally, around maps "of finite type" in a
suitable sense, which sometimes may translate into: factorizable as
a composition $f = g\circ i$ , with g proper and i an immersion. This is about
stability of "constructibility" \unsure{or|of} "finiteness" conditions for (comple-%
xes\xout{\unsure{s}} of) sheaves of k-modules\hand{,}with respect to the standard operations
$Rf_{*\hand{}}$ , $Rf_{\hand{!}}$ , $f^{\hand{!}}$ (stability by $f^{*}$ being a tautology in any case). Such
stability of course, whenever it holds, is an important feature,
for instance for making "virtual" calculations in suitable "Grothendieck
groups" (where Euler-Poincaré type invariants may be defined). It shoul
% NOTE: the right edge of the scan cuts the end of the line after "shoul"
be recalled however that the six fundamental operations in the duality
formulaire, as well as nearly all of the formulaire itself, make a
sense (and \inserted{formulae} hold\deleted{s} true) without any finiteness conditions on the complexes
of sheaves we work with, except just \deleted{\xout{restrictions} bounds} \inserted{boundedness conditions} on the degrees
of those complexes.

\margin{23}
